% history-id: model
\section{Model i warunki początkowe}

Rozpatrujemy punkt materialny wyrzucony z poziomu gruntu z prędkością
początkową \(v_0\) pod kątem \(\alpha\) do poziomu. Przyjmujemy układ
współrzędnych, w którym oś \(x\) jest pozioma, a oś \(y\) pionowa i skierowana
ku górze.

W modelu przyjmujemy:
\begin{itemize}
  \item stałe przyspieszenie grawitacyjne o wartości \(g\),
  \item brak oporu powietrza,
  \item brak innych sił działających na ciało,
  \item położenie początkowe \(\vect{r}(0)=(0,0)\).
\end{itemize}

Prędkość początkową rozkładamy na składowe
\[
v_{0x}=v_0\cos\alpha,
\qquad
v_{0y}=v_0\sin\alpha.
\]

Dla kąta \(0<\alpha<90^\circ\) obie składowe początkowe są dodatnie.
Przyspieszenie ma postać
\[
\vect{a}=(0,-g).
\]
